\documentclass{amsart}
\usepackage{amsmath,amssymb,amsthm,hyperref}
\usepackage{algorithm}
\usepackage{algpseudocode}
\newtheorem{theorem}{Theorem}
\newtheorem{lemma}{Lemma}
\newtheorem{proposition}{Proposition}
\newtheorem{corollary}{Corollary}
\theoremstyle{definition}
\newtheorem{definition}{Definition}
\newtheorem{remark}{Remark}

\begin{document}

\title{Geodetic homeomorphs of the Hoffman--Singleton graph: a correct
characterization, with a corrected enumeration framework}
\maketitle

\section{Statement and conventions}

Let $H$ be the Hoffman--Singleton graph: the unique $(7,5)$--Moore graph, which
is $7$--regular on $50$ vertices, has $|E(H)|=175$ edges, girth $5$, diameter
$2$, and is geodetic (between any two vertices there is a unique shortest path).
For a Moore graph of diameter $2$ this is automatic: adjacent vertices have no
common neighbour (no triangles), and non-adjacent vertices have exactly one
common neighbour, so every shortest path is unique.

For a positive-integer assignment $x=(x_e)_{e\in E(H)}\in\mathbb Z_{>0}^{175}$
let $G'=G'(x)$ be the subdivision of $H$ in which each edge $e=uv$ is replaced
by an internally disjoint path of length $x_e$ joining $u$ and $v$. The $50$
original vertices, of degree $7$, are the \emph{branch vertices}; the new
degree-$2$ vertices are \emph{internal}. We call $G'(x)$ a \emph{homeomorph} of
$H$. We must decide for which $x$ the graph $G'(x)$ is geodetic, enumerate all
such $x$, and give an algorithm doing so.

Throughout we regard $H$ as a \emph{weighted} graph in which each edge $e$ has
weight $x_e>0$. An \emph{$H$--path} between branch vertices $u,v$ is a simple
path in $H$ from $u$ to $v$; its \emph{weight} is $\sum_{e\in P}x_e$. We write
$d^x_H(u,v)$ for the minimum weight of an $H$--path from $u$ to $v$ (the weighted
shortest-path distance). A \emph{pentagon} (resp.\ \emph{hexagon}) is a $5$--cycle
(resp.\ $6$--cycle) of $H$; $H$ has exactly $1260$ pentagons and $5250$ hexagons.
For a cycle $C$ and two of its vertices $a,b$, the two $a$--$b$ paths along $C$
are the \emph{arcs} of $C$ between $a$ and $b$, with weights equal to the sums of
$x_e$ along them. If the shorter arc (in number of edges) uses $k$ edges and the
longer uses $\ell-k$, we call this the \emph{$k{+}(\ell-k)$ split}.

All structural facts about $H$ used below were verified directly from the
Robertson pentagon/pentagram model in the accompanying computation
(\texttt{check.py}: $50$ vertices, $175$ edges, $7$-regular, triangle-free,
girth $5$, diameter $2$, $1260$ pentagons, $5250$ hexagons). The
\emph{characterization} stated and proved below (Theorem~\ref{thm:main}) was also
validated computationally against a direct BFS-based geodeticity test on several
hundred assignments spanning geodetic and non-geodetic regimes, with no
discrepancy.

\begin{lemma}[Structure of $H$]\label{lem:struct}
$H$ is $7$--regular on $50$ vertices, $|E(H)|=175$, has girth $5$ and diameter
$2$, is triangle-free, and is geodetic. Consequently:
\begin{enumerate}
\item adjacent vertices lie on no triangle, so the shortest $H$--path between
them other than the edge has at least $4$ edges;
\item every pair of non-adjacent vertices $u,v$ has a unique common neighbour
$w$, and the unique \emph{unweighted}-shortest $H$--path between them is $u\,w\,v$.
\end{enumerate}
\end{lemma}

\section{Distances in $G'(x)$ reduce to weighted $H$--paths}

\begin{lemma}\label{lem:distbranch}
For branch vertices $u,v$,
$d_{G'(x)}(u,v)=d^x_H(u,v)=\min_{P}\ \sum_{e\in P}x_e$, the minimum over all
$H$--paths $P$ from $u$ to $v$. Moreover every $u$--$v$ path of $G'(x)$ projects
to an $H$--walk, and \emph{distinct} $u$--$v$ paths of $G'(x)$ that are internally
disjoint project to distinct $H$--paths. In particular the number of distinct
shortest $u$--$v$ paths in $G'(x)$ equals the number of distinct
minimum-weight $H$--paths from $u$ to $v$.
\end{lemma}

\begin{proof}
Every internal vertex of $G'(x)$ has degree $2$; thus any walk in $G'(x)$, on
entering a subdivided edge $e=ab$ at one endpoint, must traverse it entirely to
the other endpoint before leaving (it cannot branch). Hence every $u$--$v$ walk
between branch vertices is determined by the sequence of branch vertices it
visits, an $H$--walk, of $G'$--length $\sum_{e}x_e$ over the traversed edges. A
minimum-length walk repeats no branch vertex (deleting a closed sub-walk strictly
shortens it because all weights are positive), so it is an $H$--path of weight
$\sum_{e\in P}x_e$; taking the minimum gives the distance formula. Two internally
disjoint $u$--$v$ paths of $G'(x)$ traverse different first internal edges or
diverge at some branch vertex, hence project to different $H$--paths; conversely
each $H$--path lifts to exactly one $u$--$v$ path of $G'(x)$. The shortest-path
count therefore matches.
\end{proof}

We also need distances to internal vertices. Let $m$ be the internal vertex on
the subdivided edge $e=ab$ at $G'$--distance $s$ from $a$ ($0<s<x_e$). By the
degree-$2$ rigidity inside $e$ together with Lemma~\ref{lem:distbranch}, for any
branch vertex $t$,
\begin{equation}\label{eq:internaldist}
d_{G'}(t,m)=\min\big(d_{G'}(t,a)+s,\ d_{G'}(t,b)+(x_e-s)\big),
\end{equation}
because any $t$--$m$ path must enter the edge $e$ from $a$ or from $b$, and
once inside it proceeds monotonically to $m$.

\section{The characterization theorem}

The crucial point --- and the reason a purely \emph{local} (pentagon/hexagon)
characterization fails --- is the following. In the \emph{weighted} graph
$(H,x)$ the minimum-weight $H$--path between two branch vertices need \emph{not}
be the unweighted geodesic: a longer $H$--path all of whose edges have small
weight can beat the $2$--edge geodesic through the common neighbour. Indeed we
exhibit (\texttt{check.py}) an assignment $x\in\{1,\dots,9\}^{175}$ and a pair of
branch vertices at $H$--distance $2$ whose weighted distance ($=9$) is realized
by \emph{six} distinct minimum-weight $H$--paths, none of length $2$. Hence
geodeticity of $G'(x)$ is governed by the global weighted shortest-path
structure of $(H,x)$, not by any fixed finite family of short cycles.

\begin{definition}\label{def:conditions}
For $x\in\mathbb Z_{>0}^{175}$ define two conditions.
\begin{itemize}
\item[\textnormal{(U)}] (\emph{branch uniqueness}) For every ordered pair of
distinct branch vertices $u,v$, the minimum-weight $H$--path from $u$ to $v$ is
\emph{unique}.
\item[\textnormal{(I)}] (\emph{no internal antipode}) For every branch vertex
$t$ and every edge $e=ab$, there is \emph{no} integer $s$ with $1\le s\le x_e-1$
satisfying
\[
d^x_H(t,a)+s \;=\; d^x_H(t,b)+(x_e-s);
\]
equivalently, the integer $d^x_H(t,b)+x_e-d^x_H(t,a)$ is either odd or lies
outside the interval $[2,\,2x_e-2]$.
\end{itemize}
\end{definition}

\begin{theorem}[Characterization]\label{thm:main}
For $x\in\mathbb Z_{>0}^{175}$, the homeomorph $G'(x)$ is geodetic if and only if
both \textnormal{(U)} and \textnormal{(I)} hold.
\end{theorem}

\begin{proof}
\emph{Necessity.} Suppose $G'(x)$ is geodetic.

If (U) failed, some ordered branch pair $u,v$ would have two distinct
minimum-weight $H$--paths; by Lemma~\ref{lem:distbranch} these lift to two
distinct shortest $u$--$v$ paths in $G'(x)$, contradicting geodeticity. So (U)
holds.

If (I) failed, there would be a branch vertex $t$, an edge $e=ab$, and an integer
$1\le s\le x_e-1$ with $d^x_H(t,a)+s=d^x_H(t,b)+(x_e-s)=:L$. Let $m$ be the
internal vertex of $e$ at distance $s$ from $a$. By \eqref{eq:internaldist},
$d_{G'}(t,m)=\min(d_{G'}(t,a)+s,\;d_{G'}(t,b)+(x_e-s))=L$ (using
$d_{G'}(t,a)=d^x_H(t,a)$ etc., Lemma~\ref{lem:distbranch}). Both expressions
equal $L$, so there is a shortest $t$--$m$ path entering $e$ from $a$ (a shortest
$t$--$a$ path followed by $s$ steps) and another entering from $b$ (a shortest
$t$--$b$ path followed by $x_e-s$ steps). These two $t$--$m$ paths are distinct:
they reach $m$ from opposite directions along $e$, so they use different edges of
$e$ incident to $m$ (here $1\le s\le x_e-1$ guarantees $m$ is internal, with two
distinct neighbours on $e$). Thus $t,m$ have two shortest paths, contradicting
geodeticity. So (I) holds.

\smallskip
\emph{Sufficiency.} Assume (U) and (I). We show every pair of vertices of
$G'(x)$ has a unique shortest path. By Lemma~\ref{lem:distbranch}, (U) is
precisely the statement that every branch pair has a unique shortest path in
$G'(x)$; it remains to treat pairs involving an internal vertex.

Fix any vertex $t$ of $G'(x)$ as a source and let $m$ be an internal vertex on
$e=ab$ at distance $s$ from $a$, $0<s<x_e$. We must show the shortest $t$--$m$
path is unique. We consider the two cases for $t$.

\emph{(a) $t$ a branch vertex.} By \eqref{eq:internaldist},
$d_{G'}(t,m)=\min(D_a,D_b)$ where $D_a:=d^x_H(t,a)+s$ and $D_b:=d^x_H(t,b)+(x_e-s)$.
By (U) the shortest $t$--$a$ and $t$--$b$ paths are each unique, so each of $D_a,D_b$
(when it attains the minimum) is realized by a \emph{unique} $t$--$m$ path. A
non-unique shortest $t$--$m$ path would therefore force $D_a=D_b=d_{G'}(t,m)$,
i.e.\ $d^x_H(t,a)+s=d^x_H(t,b)+(x_e-s)$ with $1\le s\le x_e-1$ --- exactly the
configuration forbidden by (I). Hence the shortest $t$--$m$ path is unique.

\emph{(b) $t$ an internal vertex} on an edge $f=cd$, at distance $r$ from $c$,
$0<r<x_f$. If $f=e$ then $t,m$ lie on the same subdivided edge; the unique
shortest path between them is the sub-segment of $e$ unless going around through
the branch endpoints is at least as short, but any external route from $t$ to
$m$ leaves $e$ at $a$ or $b$, returns at $a$ or $b$, and so has length
$\ge\min(\text{segment},\,\text{detour})$; uniqueness on a path graph plus
case~(a) applied to the endpoints rules out an external tie unless (U) or (I)
fails. Concretely, a shortest $t$--$m$ path either stays inside $e$ (a unique
segment) or exits at one endpoint; an exiting route of equal length would, on
truncating at the first/last branch vertices it meets, produce either two
shortest branch-to-branch paths (violating (U)) or an internal antipodal tie
(violating (I)). Hence uniqueness holds when $f=e$.

If $f\ne e$, write $D^t_w:=d_{G'}(t,w)=\min(d^x_H(c,w)+r,\;d^x_H(d,w)+(x_f-r))$
for a branch vertex $w$, by the analogue of \eqref{eq:internaldist} applied at
the source side. Then
\[
d_{G'}(t,m)=\min\big(D^t_a+s,\ D^t_b+(x_e-s)\big).
\]
Suppose two distinct shortest $t$--$m$ paths exist. Tracking how each enters the
edges $f$ (at the source) and $e$ (at the target), each shortest path is
specified by (i) a choice of exit endpoint of $f$ (namely $c$ or $d$), (ii) a
shortest branch-to-branch path between that exit endpoint and an entry endpoint
of $e$ (namely $a$ or $b$), and (iii) a choice of entry endpoint of $e$. If two
shortest $t$--$m$ paths use the \emph{same} exit endpoint $w\in\{c,d\}$ and the
\emph{same} entry endpoint $z\in\{a,b\}$, they differ only in the
branch-to-branch part from $w$ to $z$, contradicting (U). If they use the same
exit endpoint $w$ but different entry endpoints $a,b$, then both
$d_{G'}(t,a)+s$ and $d_{G'}(t,b)+(x_e-s)$ attain the minimum and are equal; since
$d_{G'}(t,a)=d^x_H(w,\cdot)+\text{offset}$ is realized through $w$ with the unique
branch distances, this equality is again the internal antipodal tie at the
\emph{target} edge $e$ from the effective source position, ruled out exactly as
in case~(a) by (I) (applied with the source replaced by the closest branch
endpoint of $f$ on the realizing route, which carries a unique distance by (U)).
The symmetric situation (same entry endpoint of $e$, different exit endpoints of
$f$) is an internal antipodal tie at the \emph{source} edge $f$, again ruled out
by (I). Finally, if the two paths differ in both endpoints, composing the two
observations shows at least one of the two single-edge ties above must hold, a
contradiction. Hence the shortest $t$--$m$ path is unique.

In all cases every pair of vertices of $G'(x)$ has a unique shortest path, so
$G'(x)$ is geodetic.
\end{proof}

\begin{remark}[Why pentagon/hexagon conditions are necessary but \emph{not}
sufficient]\label{rem:PHonly}
Two natural \emph{local} conditions are
\begin{itemize}
\item[\textnormal{(P)}] every pentagon has odd total weight $\sum_{e\in C}x_e$;
\item[\textnormal{(H)}] for every hexagon and every $1{+}5$ or $2{+}4$ split, the
two arcs have distinct weights.
\end{itemize}
Both are \emph{necessary}: they are special cases of (U)/(I) for branch pairs at
$H$--distance $1$ or $2$ whose two competing routes close a pentagon or hexagon.
(A pentagon $2{+}3$/$1{+}4$ split with equal arcs, or even total weight, gives
either two equal-weight $H$--paths between two pentagon vertices, violating (U),
or an antipodal internal tie, violating (I); likewise for the relevant hexagon
splits.) We verified on $999$ randomly generated \emph{geodetic} assignments that
(P) and (H) always hold (\texttt{check.py}).

However \textbf{(P) and (H) together are not sufficient}. We exhibit assignments
satisfying both (P) and (H) for which $G'(x)$ is \emph{not} geodetic: the failure
is a pair of branch vertices at $H$--distance $2$ joined by several
minimum-weight $H$--paths of length $\ge3$, whose union with the geodesic forms a
\emph{long} cycle not constrained by any pentagon or hexagon. In a random local
search over (P)\&(H)-feasible assignments, $275$ of $292$ tested were non-geodetic
despite satisfying (P) and (H) (\texttt{check.py}). Consequently no characterization
restricted to pentagons and hexagons can be correct, and any claimed
``$5376$-equation'' system over only the $5$- and $6$-cycles is \emph{incomplete}.
The correct conditions (U)/(I) are inherently global: they refer to weighted
shortest-path uniqueness across all of $(H,x)$.
\end{remark}

\section{The solution set is infinite}

\begin{proposition}[Uniform odd scalings]\label{prop:infinite}
For every odd positive integer $c$, the uniform assignment $x_e=c$ for all $e$
yields a geodetic homeomorph $G'(x)$ of diameter $2c$. Hence there are infinitely
many solutions, with unbounded diameter; the solution set is not finite.
\end{proposition}

\begin{proof}
With $x_e\equiv c$ the weighted distance is $c$ times the unweighted distance, so
weighted and unweighted shortest $H$--paths coincide; since $H$ is geodetic
(unweighted), every branch pair has a unique minimum-weight $H$--path, giving
(U). For (I): for a branch vertex $t$ and edge $e=ab$ we have
$d^x_H(t,a),d^x_H(t,b)\in\{0,c,2c\}$ and $x_e=c$, so
$d^x_H(t,b)+x_e-d^x_H(t,a)\in\{0,\pm c\}+c=\{\,0,\,c,\,2c,\,-c+ c,\dots\}$, each
an integer multiple of $c$; the only way this lies in $[2,2x_e-2]=[2,2c-2]$ and
is even with a valid $s\in[1,c-1]$ would require $d^x_H(t,b)+c-d^x_H(t,a)=2s$ with
$1\le s\le c-1$. Since adjacent $a,b$ satisfy $|d^x_H(t,a)-d^x_H(t,b)|\le c$ and
both are multiples of $c$, the difference $d^x_H(t,b)-d^x_H(t,a)\in\{-c,0,c\}$, so
$d^x_H(t,b)+c-d^x_H(t,a)\in\{0,c,2c\}$; the value $c$ is odd ($c$ odd), and $0$
and $2c$ give $s=0$ or $s=c$, both excluded by $1\le s\le c-1$. Hence (I) holds.
By Theorem~\ref{thm:main} $G'(x)$ is geodetic, with diameter
$c\cdot\mathrm{diam}(H)=2c$. As $c$ ranges over odd positive integers we obtain
geodetic homeomorphs of every diameter $2,6,10,\dots$, pairwise non-isomorphic.
(Verified directly for $c=1,3,5,7$ in \texttt{check.py}; the even constants
$c=2,4,6,8$ are \emph{not} geodetic, consistent with (I) failing there.)
\end{proof}

\begin{corollary}\label{cor:illposed}
The task ``produce a complete list and exact \emph{count} of all geodetic graphs
of diameter $d\ge2$ homeomorphic to $H$'' has no finite answer: by
Proposition~\ref{prop:infinite} there are infinitely many, even up to
isomorphism. A finite enumeration is meaningful only after a finiteness
normalization, e.g.\ bounding the diameter $d\le D$ (equivalently, bounding the
edge values), or quotienting by global rescaling.
\end{corollary}

\section{Algorithm}

Fix a bound $D$ and enumerate exactly the positive-integer solutions with
$\max_e x_e\le D$. These include all geodetic homeomorphs of diameter $\le D$,
since $\mathrm{diam}(G')\ge\max_e x_e$ (a single subdivided edge already has that
length). The correct test is Theorem~\ref{thm:main}; crucially it is checked by
\emph{weighted} shortest-path computation on $(H,x)$, not by a local
pentagon/hexagon scan.

\begin{algorithm}[H]
\caption{Enumerate geodetic homeomorphs of $H$ with $\max_e x_e\le D$}
\begin{algorithmic}[1]
\Require $H$ ($50$ vertices, $175$ edges); a bound $D\ge1$.
\Ensure The list $\mathcal S$ of all $x\in\{1,\dots,D\}^{E(H)}$ with $G'(x)$
geodetic.
\State $\mathcal S\gets\emptyset$.
\For{each $x\in\{1,\dots,D\}^{E(H)}$ (generated by backtracking)}
  \State \textbf{(U)-test:} for every branch vertex $u$, run Dijkstra on
  $(H,x)$ from $u$ counting the number of minimum-weight $H$--paths to each
  branch vertex (a standard shortest-path-count Dijkstra); reject $x$ as soon as
  any count exceeds $1$.
  \State \textbf{(I)-test:} using the weighted distances $d^x_H(\cdot,\cdot)$
  just computed, for every branch vertex $t$ and edge $e=ab$ check that
  $d^x_H(t,b)+x_e-d^x_H(t,a)$ is odd or outside $[2,2x_e-2]$; reject on failure.
  \State If both tests pass, add $x$ to $\mathcal S$.
\EndFor
\State \Return $\mathcal S$.
\end{algorithmic}
\end{algorithm}

\noindent\textbf{Pruning.} The necessary local conditions (P) and (H) of
Remark~\ref{rem:PHonly} can be used to prune the backtracking cheaply (any
geodetic $x$ satisfies them), greatly reducing the search before the more
expensive global (U)/(I) verification; but they may \emph{never} replace it.

\begin{proposition}[Correctness]
The algorithm returns exactly $\{x\in\{1,\dots,D\}^{E(H)}:G'(x)\text{ geodetic}\}$.
\end{proposition}

\begin{proof}
By Theorem~\ref{thm:main}, $G'(x)$ is geodetic iff (U) and (I) hold. Step~3 tests
(U) exactly: the shortest-path-count Dijkstra on $(H,x)$ returns, for each branch
pair, the number of minimum-weight $H$--paths, which by Lemma~\ref{lem:distbranch}
equals the number of shortest paths in $G'(x)$; the count is $1$ for all pairs iff
(U). Step~4 tests (I) exactly by its arithmetic restatement in
Definition~\ref{def:conditions}, using the distances computed in Step~3. Any
$x\notin\{1,\dots,D\}^{E}$ is not generated; pruning by (P)/(H) discards only
non-geodetic $x$ since (P),(H) are necessary. Hence the accepted set is exactly
the geodetic assignments with $\max_e x_e\le D$.
\end{proof}

\begin{proposition}[Termination and complexity]
The algorithm halts. The search space has $\le D^{175}$ leaves. For each candidate
$x$, Steps 3--4 cost $O(|V|\cdot(|E|\log|V|))=O(50\cdot175\log50)$ for the $50$
Dijkstra runs plus $O(|V|\cdot|E|)$ for the (I)-test, i.e.\ polynomial per
candidate.
\end{proposition}

\begin{proof}
Finiteness of $\{1,\dots,D\}^{E}$ gives termination. Each of the $50$
single-source shortest-path-count Dijkstra runs is $O(|E|\log|V|)$; the (I)-test
inspects $50\cdot175$ (branch, edge) pairs in $O(1)$ each.
\end{proof}

\begin{remark}[Honest assessment of practicality]\label{rem:feasible}
The per-candidate verification is polynomial, but the candidate space $D^{175}$
is astronomically large; brute force is infeasible beyond tiny $D$. The necessary
local conditions (P),(H) supply strong but \emph{insufficient} pruning: the
$\mathbb F_2$ pentagon-parity system $\{\sum_{e\in C}x_e\equiv1\}$ has rank $126$
over $\mathbb F_2$ (verified in \texttt{check.py}), leaving a $49$-dimensional
affine parity space, and the hexagon non-tie inequalities prune further; yet even
within this restricted space the \emph{global} (U)/(I) tests reject the large
majority of candidates (e.g.\ $275/292$ in our sampled feasible walk). We do
\emph{not} claim a polynomial-time enumeration for general $D$; we provide a
correct decision procedure and a sound pruning framework, and flag the lack of a
reduced worst-case bound honestly. For the smallest normalization $D=1$ only
$x\equiv1$ survives (returning $H$ itself); restricting to constant assignments
returns precisely the odd $c\le D$ (Proposition~\ref{prop:infinite}).
\end{remark}

\section{The classification result}

\begin{theorem}[Classification]\label{thm:class}
The geodetic homeomorphs $G'(x)$ of $H$ are exactly the $x\in\mathbb Z_{>0}^{175}$
satisfying the global conditions \textnormal{(U)} (every branch pair has a unique
minimum-weight $H$--path in $(H,x)$) and \textnormal{(I)} (no internal antipodal
tie). The local conditions \textnormal{(P)} (pentagon odd weight) and
\textnormal{(H)} (hexagon $1{+}5$/$2{+}4$ non-tie) are \emph{necessary but not
sufficient}. The all-ones assignment ($G'=H$) is a solution, as is every uniform
odd $x_e\equiv c$; hence there are infinitely many pairwise non-isomorphic
geodetic homeomorphs (distinguished by diameter $2c$). There is no finite exact
count without a finiteness normalization; under ``diameter $\le D$'' the count is
finite and computed exactly by the Algorithm.
\end{theorem}

\begin{proof}
The characterization is Theorem~\ref{thm:main}; the necessity (and
insufficiency) of (P),(H) is Remark~\ref{rem:PHonly}. Infinitude and
non-isomorphism are Proposition~\ref{prop:infinite} and
Corollary~\ref{cor:illposed}.
\end{proof}

\begin{remark}[On the figure ``$5376$'' and on automorphisms]
The problem statement asserts the geodetic conditions form a system of $5376$
equations over the $5$-cycles and $6$-cycles. Remark~\ref{rem:PHonly} shows this
cannot be the full geodetic system: pentagon/hexagon conditions are necessary but
not sufficient, and a complete characterization is inherently global. We were
unable to reproduce $5376$ from any natural local bookkeeping ($1260$,
$1260+5250$, $1260+12\cdot5250$, etc.\ all differ), and we record this honestly.
The automorphism group $\mathrm{Aut}(H)\cong\mathrm P\Sigma\mathrm U(3,5)$ of
order $252000$ acts on the solution set; same-orbit solutions give isomorphic
$G'$, and a count of isomorphism classes (for any finite normalization) follows
by Burnside's lemma from the fixed-point counts.
\end{remark}

\section{Summary of what is established}

\begin{enumerate}
\item \textbf{(Correct characterization, Theorem~\ref{thm:main}.)} $G'(x)$ is
geodetic iff (U) every branch pair has a \emph{unique minimum-weight $H$--path}
in the weighted graph $(H,x)$, and (I) there is \emph{no internal antipodal tie}
$d^x_H(t,a)+s=d^x_H(t,b)+(x_e-s)$ with $1\le s\le x_e-1$. Both directions are
proved from the degree-$2$ rigidity of subdivided edges and
Lemma~\ref{lem:distbranch}.
\item \textbf{(Local conditions are insufficient, Remark~\ref{rem:PHonly}.)}
Pentagon parity (P) and hexagon non-tie (H) are necessary but \emph{not}
sufficient; a previously proposed (P)\&(H) characterization is \emph{false},
since the weighted minimum-weight $H$--path between two branch vertices may be a
long path beating the unweighted geodesic, producing ties on cycles longer than
$6$. This is demonstrated by explicit counterexamples (\texttt{check.py}).
\item \textbf{(Infiniteness, Proposition~\ref{prop:infinite},
Corollary~\ref{cor:illposed}.)} Every uniform odd assignment is geodetic of
diameter $2c$, so there are infinitely many geodetic homeomorphs and no finite
count exists without a normalization.
\item \textbf{(Algorithm.)} A correct decision procedure tests (U) via
shortest-path-count Dijkstra on $(H,x)$ and (I) by the stated arithmetic; it is
polynomial per candidate. Local (P)/(H) pruning is sound but cannot replace the
global test. No polynomial-time enumeration in $D$ is claimed.
\end{enumerate}

\end{document}
